% !PRINTED 494
\weberpage{494}

%<<CONT
imprimitive Zahlen des K\"orpers unterschieden. Eine symmetrische Function
der conjugirten Zahlen ist eine rationale Zahl. Unter diesen symmetrischen
Functionen sind zwei von besonderer Wichtigkeit, die Summe und das Product,
von denen die erste die \gesperrt{Spur}, die zweite die \gesperrt{Norm} von
$\omega$ genannt wird. Man bezeichnet diese beiden Zahlen durch $S(\omega)$
und $N(\omega)$.

Hierbei werden, wenn unter den conjugirten Zahlen dieselben Zahlen mehrfach
vorkommen, diese gleichen Zahlen so oft in die Summe oder das Product
aufgenommen, als der Grad ihrer H\"aufigkeit angiebt.

Da sich in jedem solchen Zahlk\"orper die vier fundamentalen
Rechenoperationen ebenso wie im K\"orper der rationalen Zahlen ausf\"uhren
lassen, so kann man auch die Frage aufwerfen, inwieweit sich die aus der
Theorie der rationalen Zahlen bekannten arithmetischen Grundgesetze in einem
beliebigen algebraischen Zahlk\"orper bew\"ahren. Es handelt sich hierbei in
erster Linie um die Zerlegung der ganzen Zahlen in ihre \gesperrt{Primfactoren}.

Da diese Zerlegung mit den Zahlen des algebraischen K\"orpers selbst im
Allgemeinen \gesperrt{nicht} gelingt, so ist eine Erweiterung des
Rechenmateriales n\"othig, um die einfachen Gesetze wieder herzustellen, und
eine solche Erweiterung ist in verschiedenem Sinne m\"oglich. Es m\"ussen
sich aber diese verschiedenen Erweiterungen auf einander zur\"uckf\"uhren,
oder, genauer gesagt, in eine eindeutige Beziehung zu einander setzen
lassen.

\gesperrt{Kummer} hat zuerst f\"ur die aus Einheitswurzeln gebildeten
algebraischen Zahlen (die Kreistheilungszahlen) das grosse Problem durch die
Sch\"opfung der \gesperrt{idealen Zahlen}\textsuperscript{1)} gel\"ost. Eine
andere ganz allgemeine, keiner Ausnahme unterworfene L\"osung hat die
Aufgabe durch \gesperrt{Dedekind} gefunden, der als die einfachsten Elemente
der Rechnung die von ihm so genannten \gesperrt{Ideale}\textsuperscript{2)}

\vspace{1em}\hrule\vspace{0.5em}

\noindent\textsuperscript{1)} \gesperrt{Kummer}, Theorie der idealen
Primfactoren der complexen Zahlen etc. \gesperrt{Crelle}'s Journal, Bd.\,35, 1846,
Bd.\,40, 1850. Abhandlungen der Berliner Akademie 1856. Sur la th\'eorie des
nombres complexes compos\'es de racines de l'unit\'e et de nombres entiers.
\gesperrt{Liouville}'s Journal, Bd.\,16, 1851.

\noindent\textsuperscript{2)} \gesperrt{Dedekind}, in dem letzten Supplement
der 2., 3. und 4. Auflage von \gesperrt{Dirichlet's} Vorlesungen \"uber
Zahlentheorie (Braunschweig 1871, 1879, 1894). Zu vergleichen ist auch: Sur
la th\'eorie des nombres entiers alg\'ebriques im Bulletin von
\gesperrt{Darboux} und \gesperrt{Ho\"uel} (1\textsuperscript{re} s\'er. XI,
1877). Ueber den Zusammenhang zwischen der Theorie der Ideale und der
h\"oheren Congruenzen (Abhandlungen der Ges. der Wissensch. in G\"ottingen,
Bd.\,23,
%CONT>>
